\scanpage{38}
Le sous-groupe \inserted{$H\subset\underline{\operatorname{GP}}(V)$} bien sûr ne dépend que du sous-Module $\unsure{\mathcal O_S|\mathcal O_S1}+\mathcal O_Su$ de $\underline{\operatorname{End}}(V)$, ou ce qui revient au $=$, de son image dans $\underline{\operatorname{End}}(V)/\underline{\mathcal O}_S.1=\underline{\operatorname{Lie}}\,\underline{\operatorname{GP}}(V)$ laquelle image est un sous-fibré inversible qui n'est autre que $\operatorname{Lie}H$. Ceci précise relation entre sous-fibrés vectoriels \inserted{$\mathfrak h$} de rang 1 de $\mathfrak g=\underline{\operatorname{GP}}(V)=G$ et sous-gpes lisses \inserted{connexes et comm. $H$} de \deleted{rang} \inserted{dim rel.} 1 de $\unsure{G|\mathfrak g}$.

Est-il vrai que $H$ est le centralisateur de $\mathfrak h$ \inserted{$=\operatorname{Lie}\unsure{\mathfrak h|H}$}, ou de tte section de $H$ partout $\ne$ section unité, ? à fortiori de $H$ lui-$=$ ?
\[
\begin{aligned}
ghg^{-1}\deleted{\unsure{\ne}}&=\lambda h
& \text{sps }\ \deleted{\unsure{u\ \mathrm{univ.}|\mathrm{unipotent}}}\,\inserted{\det u=-1}\\
\det h&=\lambda^2\det h & \boxed{\lambda^2=1}\\
\deleted{\operatorname{Tr}u}&\deleted{=\lambda\operatorname{Tr}u}
\end{aligned}
\]
\deleted{en car. $\ne2$ on déduit $\lambda=\pm1$ loc$^{\mathrm{t}}$, le cas $+1$ déjà traité, reste \inserted{alors} le cas}
\[
 \deleted{gug^{-1}=-u}
\]
\deleted{Ceci \unsure{est un} n'est possible que si \deleted{les valeurs} \inserted{$\operatorname{Tr}u=0$, donc} (avec $\det u=-1$) est donc \inserted{$u$ une réflexion} \deleted{propres de $u$ sont de la forme} $\{1,-1\}$ (donc distinctes, \deleted{\unsure{et si elles sont}} \inserted{puisqu'elles sont} non nulles), $=$ est conjugué \unsure{à une|une} réflexion dans $G$ donc $u^2=1$} Si on est en car rés. 2,
% archive sheet number: 37
